The main goal of the \textit{Best Bike Paths (BBP)} system is to provide cyclists with a mobile application that allows them to record their activities and share information about bike routes, while also fostering the creation of an active and collaborative community.
The application enables registered users to save their rides, view detailed statistics such as total distance covered or average speed, and contribute to updating the path inventory by reporting the condition of routes or the presence of potential obstacles.
The information entered by users can be made public and is then aggregated by the system, which combines data from multiple sources to offer a comprehensive and up-to-date representation of the paths’ conditions.
Unregistered users can still take advantage of the consultation features by searching for routes between two points and viewing the available options on a map, ranked according to a score that considers both the quality and effectiveness of the routes.
The system also interacts with an external meteorological service, which provides weather data to enrich the information related to recorded activities.

\subsection{World Phenomena}
\begin{enumerate}[label=WP\arabic*.]
    \item Cyclists travel through real-world routes by following predefined paths.
    \item Cyclists perceive the quality of the path.
    \item Cyclists search for paths between two real-world points.
    \item The physical condition of paths can vary over time.
    \item Mobile devices can detect data about the cyclist in the real world.
    \item Real-world weather conditions are detected by meteorological services.
    \item There are bike lanes and roads where speed limits or the absence of cars make them suitable for bicycles.
\end{enumerate}

\subsection{Shared Phenomena}
\textbf{User-System}
\begin{enumerate}[label=SP\arabic*.]
    \item Users can register within the application.
    \item The user can record and store their cycling activities.
    \item The user can choose whether to publish route-related data.
    \item The system can provide the user with statistics about paths.
    \item The user can manually enter information about paths.
    \item The system can automatically record information, and the user must confirm it.
    \item The user searches for paths based on a starting point and an ending point.
\end{enumerate}
\textbf{System-Device}
\begin{enumerate}[label=SP\arabic*., resume]
    \item The system receives data from mobile devices to track routes and detect obstacles.
    \item The system sends requests to an external service and receives weather data.
\end{enumerate}